\documentclass[11pt]{article}
\usepackage[margin=1in]{geometry}
\usepackage{amsmath,amssymb,amsthm,mathtools}
\usepackage{bm}
\usepackage{hyperref}
\usepackage{enumitem}
\usepackage{booktabs}
\usepackage{array}
\usepackage{longtable}
\usepackage{graphicx}
\usepackage{setspace}
\usepackage{microtype}
\usepackage[T1]{fontenc}
\usepackage[utf8]{inputenc}

\title{Relational Actualism: Predictions, Departures, and Tests}
\author{Joshua Sandeman}
\date{\today}

\begin{document}
\maketitle

\begin{abstract}
Relational Actualism (RA) is not only an ontological and mathematical reconstruction of physics. It is also an exposed empirical program. If the world is fundamentally a growing causal ledger of irreversible actualization events, if geometry is the dense-regime order of that ledger, if matter is stable local self-reproduction, and if persistence arises through filtered exit geometry, then the framework should produce characteristic empirical consequences. This paper collects those consequences.

The purpose of the paper is not to reargue the kernel ontology or the full bridge program. Those tasks belong to the companion theory-definition and bridge papers. The present paper asks a different question: what does implemented $d=4$ Relational Actualism predict, where does it depart from standard frameworks, and what would count as empirical failure?

The paper develops a prediction architecture in three layers. First are direct discriminator predictions: empirical consequences tied closely to the actualization program itself, such as irreversibility thresholds, macroscopic record formation, and null or non-null signatures in experiments that test whether gravity alone can induce genuine quantum actualization. Second are structural departures: consequences for cosmology, dark-sector phenomenology, conservation laws, and the persistence structure of matter once the universe is modeled as a self-writing, self-filtering graph rather than as a continuum state space. Third are higher-order persistence consequences: predictions concerning life, cognition, and the distribution of complex organization, treated not as additions from outside physics but as recursive stabilization phenomena in the same graph-theoretic framework.

A central methodological claim of the paper is that RA should be judged not only by whether it reproduces known structures, but by whether it exposes itself to characteristic failure modes. The framework therefore does not present its predictions as a homogeneous block. Each prediction is traced, as far as possible, to its dependency structure: kernel commitment, closure result, implementation detail, bridge layer, and current epistemic status. Some predictions are close to the discrete substrate and relatively robust. Others are farther downstream and should be read explicitly as structural or exploratory consequences rather than theorem-level outputs.

The paper also develops a unifying theme that has become clearer in recent work: persistence is not primitive and decay is not accidental. Across scales, structures survive when cheap exit paths are filtered or blocked, and fail when they are not. This principle appears first in the particle lifetime program and is then extended cautiously into the domains of cosmology, chemistry, biology, and agency. In this sense, the predictive layer of RA is not merely a list of disconnected tests. It is the empirical face of a self-organizing and self-stabilizing universe.
\end{abstract}

\section*{Front-matter status box}
\noindent\textbf{This paper does:} collect the framework's testable predictions; separate direct discriminator predictions from broader structural departures; and trace each major prediction back to the kernel, closure, implementation, and bridge layers on which it depends.

\noindent\textbf{This paper does not:} redefine the ontology, rederive the full bridge program, or present interpretive mappings as theorem-level outputs. Those tasks belong to the companion papers.

\noindent\textbf{Status discipline:} predictions are marked in prose as one of the following: \textbf{RB} (relatively robust, close to kernel/closure), \textbf{BR} (bridge-dependent), \textbf{SE} (structural extension), or \textbf{EX} (exploratory).

\section{Introduction: Why predictions matter here}
A new framework in fundamental physics is not tested only by whether it can be made philosophically coherent or mathematically nontrivial. It is tested by whether it exposes itself to the world. Relational Actualism is no exception. If the universe is fundamentally a growing causal ledger of irreversible events, if time is graph growth rather than a background parameter, if matter is stable local self-reproduction, and if persistence depends on filtered exit geometry, then the framework should generate empirical consequences that are not generic to continuum state-based physics.

This paper collects those consequences. It is the predictive layer of the canonical RA suite. Paper~I defines the ontology, closure architecture, and implemented $d=4$ theory. Paper~II develops the strongest bridge results to known physics: gravity, finite topology closure, interaction hierarchy, and the motif-renewal lifetime program. The purpose of the present paper is different. It asks what follows empirically when one treats those results not as isolated achievements but as parts of a single process-based framework.

The predictions presented here differ significantly in epistemic status. Some sit very close to the actualization program itself. Others depend on the full bridge architecture from density-controlled exit geometry and filtered persistence. Still others are structural extensions of the same ideas into domains such as astrobiology and higher-order organization. It is therefore a mistake to read the paper as offering one homogeneous block of equally secure outputs. The framework is strongest when it states explicitly which predictions are close to the discrete substrate and which depend on bridge or extension layers.

A second organizing principle is also important. Recent work has made clearer that persistence in RA is not the default state of things. Stable structure exists because graph growth is self-filtering: renewable motifs survive, cheap exits are blocked or narrowed, and increasingly organized systems persist only to the extent that their internal organization filters dominant disruption channels. This first appears in the particle-lifetime program, but it naturally suggests a broader predictive picture. Cosmological departures, dark-sector alternatives, and higher-order closure phenomena may all be read as consequences of how graph growth stabilizes or fails to stabilize particular structures.

The paper is organized accordingly. Section~\ref{sec:taxonomy} introduces a prediction taxonomy. Section~\ref{sec:actualization-preds} presents predictions closest to the actualization criterion itself. Section~\ref{sec:gravity-cosmology} presents gravity and cosmology departures. Section~\ref{sec:matter-conservation} presents matter-sector and conservation predictions. Section~\ref{sec:higher-order} develops higher-order persistence consequences, including the causal-firewall idea and biosignature consequences. Section~\ref{sec:dependency-table} presents a dependency framework for evaluating the predictions. Section~\ref{sec:conclusion} closes.

The goal is not to maximize rhetorical reach. It is to present RA as a framework that now has enough structure to risk contact with the world. That is what a predictive paper should do.

\section{Prediction taxonomy}\label{sec:taxonomy}
\subsection{Why a taxonomy is needed}
A major weakness of many ambitious foundational programs is that they present all consequences in the same tone. Near-theory outputs, bridge-dependent extrapolations, and speculative long-range consequences appear side by side without sufficient status discipline. RA should not do this. The predictive layer is strongest when it explicitly distinguishes the proximity of a claim to the framework's core.

\subsection{Four prediction classes}
The predictions in this paper are therefore grouped into four classes.

\paragraph{RB --- Relatively robust.}
These are predictions close to kernel, closure, or implementation structure. They do not require long interpretive ladders. They are not necessarily theorem-level, but they are the hardest to dismiss as arbitrary extensions.

\paragraph{BR --- Bridge-dependent.}
These depend on one or more bridge layers developed in Paper~II, such as the gravity bridge, the topology-to-particle mapping, or the motif-renewal lifetime program.

\paragraph{SE --- Structural extension.}
These are not arbitrary speculations, but consequences of extending a bridge result into a neighboring domain, such as cosmology or chemistry.

\paragraph{EX --- Exploratory.}
These are conjectural consequences of the framework whose motivating logic is visible, but whose dependence chain is longer and whose closure status remains weaker.

\subsection{A second axis: falsifiability}
In addition to epistemic depth, predictions differ in how sharply they expose the framework to failure. Some are close to direct experimental decision. Others are presently only contrastive or pattern-based. This matters because a framework should not merely announce consequences; it should state what kinds of empirical outcomes would count against it.

\section{Actualization-level predictions}\label{sec:actualization-preds}
\subsection{Main status of this section}
\textbf{Main status:} \textbf{RB/BR}.

These predictions sit closest to the actualization program itself. They are therefore among the most load-bearing predictions in the suite.

\subsection{Irreversibility thresholds}
If actualization is fundamentally the irreversible inscription of a previously open possibility into the causal ledger, then there should exist a nontrivial threshold separating reversible possibility dynamics from record-forming actuality. The discrete BDG filter is the native expression of that threshold. A strong empirical expectation therefore follows: there should be domains in which apparently interaction-rich processes nevertheless fail to generate true record-forming actualization, and domains in which threshold-crossing transitions become sharply irreversible.

The exact experimental expression of this expectation is still an open bridge problem, but the broad prediction is robust: macroscopic irreversibility should not be fundamental everywhere and always; it should emerge where graph growth reaches the relevant threshold for record formation.

\subsection{BMV-type tests}
One of the most visible near-term discriminator classes concerns whether gravity alone can induce genuine quantum actualization. In the RA picture, if gravitational interaction remains below the relevant threshold for irreversible inscription, then entanglement-style signatures may occur without actualization. Conversely, if gravity in a sufficiently strong configuration does cross the threshold, then the result should not be mere dephasing but genuine record formation.

This is one of the most important open test arenas because it probes the actualization threshold directly rather than only its downstream continuum paraphrases. The exact sign of the prediction depends on the operational realization, but the broader expectation is clear: gravity should not be treated as a merely passive background in such experiments, because in RA the distinction between possibility and actuality is dynamical and thresholded.

\subsection{Loschmidt-echo and reversibility tests}
RA also motivates a family of reversibility-threshold predictions. If actualization is primitive and irreversible, then there should be experimentally meaningful distinctions between systems that remain within the possibility layer and systems that have crossed into genuine causal inscription. This suggests that Loschmidt-echo–type restoration experiments should show a structurally significant dependence on whether record-forming thresholds have been crossed. In the deepest form of the claim, reversibility is not merely hard in practice; it is blocked in principle once irreversible inscription has occurred.

\subsection{What would count against RA here}
RA would be pressured if evidence accumulated that apparently macroscopic, information-bearing processes could be reversed without residue even after the conditions the framework identifies with record formation had been crossed. It would also be pressured if every plausible actualization-sensitive experiment could be modeled without any threshold distinction between reversible possibility evolution and irreversible inscription. The framework survives only if the actualization distinction becomes empirically meaningful somewhere.

\section{Gravity and cosmology departures}\label{sec:gravity-cosmology}
\subsection{Main status of this section}
\textbf{Main status:} \textbf{BR/SE}.

The gravity bridge in Paper~II argues that implemented $d=4$ RA supports Einstein dynamics in the dense-regime continuum limit. The present section concerns what happens when one moves away from that regime or asks what cosmological consequences follow when the universe is interpreted as a self-writing graph rather than as a continuum metric field theory with an independently given vacuum structure.

\subsection{No fundamental $\Lambda$ term}
One of the most characteristic structural consequences of RA is that the framework does not begin with a primitive vacuum energy term that must then be explained away. The growing-ledger ontology strongly suggests a universe whose large-scale order emerges from actualization structure rather than from a bare cosmological constant inserted at the fundamental level. The strongest current form of the claim is therefore that $\Lambda$ is not a primitive term in the discrete substrate. This is not by itself a full cosmology, but it is already a major departure from standard effective-field reasoning.

\subsection{Sparse-regime departures}
If Einstein dynamics is a dense-regime continuum limit of a discrete actualization process, then one should expect departures where the graph is sufficiently sparse or weakly correlated that the continuum approximation begins to fail. Such departures are not expected to be arbitrary. They should appear where the assumptions behind the continuum closure are weakest. This motivates a class of structural alternatives to standard dark-sector explanations: some apparent anomalies may reflect sparse-regime deviations in the actualization graph rather than unseen matter fields.

\subsection{Hubble-gradient and matter-distribution effects}
RA also motivates the possibility that large-scale expansion behavior depends on graph-density and matter-distribution structure more directly than in standard homogeneous continuum cosmology. If causal correlation and stabilization history matter, then one should expect environment-sensitive expansion features or gradients. At present this remains a structural extension rather than a closed prediction, but it is a natural place where RA departs from a picture in which spacetime geometry is primary and matter merely sources it.

\subsection{Rotation-axis and nucleation-style anomalies}
A more speculative but still structurally motivated class of departures concerns preferred large-scale features seeded by early graph nucleation structure, including axis-like or Kerr-like asymmetries. These ideas should remain status-labeled as exploratory, but they are not arbitrary additions. They follow from taking the initial growth structure of the ledger seriously rather than assuming from the start that cosmological geometry is exactly featureless.

\subsection{What would count against RA here}
RA would be pressured if all well-confirmed cosmological departures from naive homogeneous GR required structures that could not plausibly be interpreted as sparse-regime or graph-growth effects. It would also be pressured if the only viable cosmological fits required reinstating a primitive vacuum term as a basic ingredient rather than a late effective summary.

\section{Matter and conservation predictions}\label{sec:matter-conservation}
\subsection{Main status of this section}
\textbf{Main status:} \textbf{BR}.

This section collects predictions tied to the particle-topology bridge, motif-renewal bridge, and the conservation structure of the implemented theory.

\subsection{Exact baryon conservation}
One of the strongest matter-sector consequences of the topology and renewal program is that baryon persistence is not an accident of approximate low-energy symmetry. It is tied to the topology-space structure of available exits. In the strongest form of the claim, exact baryon conservation follows because the relevant motif classes do not admit the required destabilizing pathways. This should be treated as a bridge-dependent prediction rather than a theorem-level fact of the kernel, but it is one of the framework's boldest and cleanest contrasts with many beyond-Standard-Model programs.

\subsection{WIMP prohibition and dark-sector austerity}
If matter is restricted to the finite stable motif universe and if persistence requires specific renewal/exit structures, then broad classes of ad hoc stable dark-sector particles are disfavored or excluded. The exact boundary of this exclusion depends on how tightly one reads the topology closure theorem and associated bridge arguments, but the general prediction is clear: the framework does not naturally support an unrestricted zoo of stable exotic particle species.

\subsection{Neutrino and light-sector constraints}
The light sector remains less closed than the hadronic resonance sector, but the general RA logic still imposes constraints. Long-lived or weakly interacting motifs must correspond to highly filtered exit geometry, and the same framework that distinguishes direct, strange, parity-blocked, and topology-protected decays in the hadron sector should eventually constrain the neutrino sector more tightly than standard parameter-rich approaches do. At present this remains bridge-level and partly exploratory.

\subsection{Strong-decay lifetime hierarchy}
Paper~II develops the direct strong-decay baseline, the hadronic density closure, the stabilization-layer hierarchy, and the pathwise exit program. The present paper treats the resulting lifetime hierarchy as a prediction family. Direct-exit motifs should be shortest-lived. Flavor-rearrangement motifs should be longer-lived. Multi-step symmetry-filtered motifs should be longer still. Topology-protected motifs should exhibit no cheap one-step exit. This hierarchy is one of the clearest examples of a predictive structure that is both novel and rooted in the discrete substrate.

\subsection{What would count against RA here}
The framework would be pressured if a broad and well-confirmed stable particle sector were discovered that could not be accommodated within any reasonable extension of the finite topology universe. It would also be pressured if proton decay were observed in a channel that the topology-space picture genuinely forbids. More generally, the bridge would be weakened if the lifetime hierarchy ceased to track the accessibility structure of exit geometry.

\section{Higher-order persistence: chemistry, life, and agency}\label{sec:higher-order}
\subsection{Main status of this section}
\textbf{Main status:} \textbf{SE/EX}.

This section should be read with more caution than the earlier ones. It is not theorem-level closure, nor even as strong as the particle-bridge layer. But it is also not disconnected speculation. It extends a principle that has already appeared in the motif-lifetime program: stable structure persists when cheap disruption channels are filtered or blocked while renewal remains possible.

\subsection{From particle stabilization to higher-order stabilization}
The strongest unifying idea is this: persistence across scales may be governed by the same general logic. At the particle level, a motif survives because direct exits are filtered, lengthened, or blocked. At higher scales, organized systems may persist for the same general reason: their internal organization filters dominant disruption channels while preserving renewable internal flow. If this is correct, then the distinction between physics and higher-order organization is not a jump from one ontology to another. It is a recursive continuation of the same stabilization principle.

\subsection{The causal firewall}
One of the most distinctive higher-order proposals in RA is the causal-firewall idea. In its strongest disciplined form, the claim is not that living systems are magically exempt from ordinary physics. It is that living systems may occupy regions of motif space whose internal organization blocks the cheapest thermal and chemical disruption pathways. On that reading, a causal firewall is not a mystical boundary but an organization-dependent exit filter.

The apt analogy suggested by recent work is that what OZI-like topology protection is in the hadron sector, high-order disruption-channel filtering may be at the molecular or biological scale. This should remain a structural extension, not a theorem claim. But it is a deep and promising one.

\subsection{Biosignature consequences}
If life corresponds to a regime of unusually strong exit-path filtering together with sustained internal renewal, then the search for life should not focus only on chemistry in the static sense. It should focus on systems whose organization appears to block dominant degradation channels while maintaining renewable flow. That suggests a richer biosignature logic than one based only on equilibrium chemical composition.

\subsection{Agency and cognition}
The same logic extends, more speculatively, to agency and cognition. Systems capable of maintaining long-lived, recursively stabilized internal states in the face of disruption may correspond to high-order motifs whose internal structure creates increasingly selective filters on destructive transitions. Again, this should remain exploratory. But it is no longer arbitrary once the lower-scale motif-renewal program is taken seriously.

\subsection{What would count against RA here}
The higher-order extension would be weakened if no meaningful empirical distinction could be drawn between systems that merely persist passively and systems whose organization clearly filters dominant disruption channels. It would also be weakened if the causal-firewall idea proved to have no empirical bite at all beyond metaphor.

\section{Prediction dependency framework}\label{sec:dependency-table}
\subsection{Why dependencies matter}
Predictions are strongest when their dependency structure is visible. A claim resting directly on the actualization criterion is not in the same epistemic position as a cosmological consequence that depends on the gravity bridge, sparse-regime departures, and a particular interpretation of early graph nucleation. This section therefore states how predictions in RA should be read.

\subsection{Dependency schema}
Every major prediction in the framework should, where possible, be traced through the following levels:
\begin{enumerate}[label=\arabic*.]
    \item \textbf{Kernel source:} which kernel commitment is relevant?
    \item \textbf{Closure source:} which closure result pins the relevant structure?
    \item \textbf{Implementation source:} what part of implemented $d=4$ RA is being used?
    \item \textbf{Bridge source:} which bridge program, if any, carries the consequence into familiar physics?
    \item \textbf{Extension layer:} is the claim direct, bridge-dependent, structural, or exploratory?
    \item \textbf{Failure mode:} what empirical outcome would genuinely count against it?
\end{enumerate}

\subsection{Illustrative examples}
\begin{longtable}{>{\raggedright\arraybackslash}p{3.2cm} >{\raggedright\arraybackslash}p{2.3cm} >{\raggedright\arraybackslash}p{2.3cm} >{\raggedright\arraybackslash}p{2.3cm} >{\raggedright\arraybackslash}p{1.5cm} >{\raggedright\arraybackslash}p{3.0cm}}
\toprule
Prediction & Kernel / closure source & Implementation / bridge source & Extension class & Status & Clear failure mode \\
\midrule
Irreversibility threshold & actuality threshold + selectivity structure & BDG filter / actualization hierarchy & direct & RB & no empirical distinction between reversible possibility and irreversible inscription \\
Strong-decay lifetime hierarchy & stable motifs + finite topology + density closure & motif-renewal / pathwise exit bridge & bridge & BR & hierarchy fails to track exit accessibility structure \\
No primitive $\Lambda$ term & event-first ontology + dense-limit bridge & gravity bridge / sparse-regime extension & structural & SE & only viable cosmology requires primitive vacuum term \\
Exact baryon conservation & finite motif universe + blocked exits & topology / renewal bridge & bridge & BR & confirmed proton decay in forbidden channel \\
Causal firewall & renewable motifs + filtered exits & high-order stabilization extension & exploratory & EX & no empirically meaningful disruption-filtering distinction in living systems \\
\bottomrule
\end{longtable}

This dependency discipline should be treated as part of the framework itself. It is not merely a presentational aid.

\section{Conclusion}\label{sec:conclusion}
This paper has one central goal: to present Relational Actualism as an empirical program rather than only as an ontological or mathematical one. The framework now has enough internal structure that this can be done with real discipline. Some of its predictions sit close to the actualization criterion and should be regarded as relatively robust. Some depend on the bridge architecture developed in Paper~II. Others extend those bridge ideas into broader domains such as cosmology, chemistry, and life. The framework is strongest when these levels remain explicit.

The most important unifying thread is that persistence is not primitive and decay is not accidental. In the direct particle sector, this appears as density-controlled exit geometry and stabilization layers. In the more filtered lifetime hierarchy, it appears as selection rules reinterpreted as internal consistency filters on topology-space exits. In higher-order domains, it appears as the possibility that increasingly organized structures persist by recursively filtering dominant disruption channels. This principle does not prove every specific extension offered here. But it does explain why the predictions are not merely a list of unrelated speculations.

The framework therefore exposes itself to the world in several distinct ways. It predicts that actualization is thresholded and empirically meaningful. It predicts that the direct strong-decay sector should be governed by a closure-backed density scale and that longer lifetimes arise through additional stabilization layers. It predicts a more austere matter sector than many extension-heavy continuum frameworks. It predicts that some apparent cosmological anomalies may reflect sparse-regime departures of a self-writing graph rather than new invisible fields. And it suggests that higher-order organization may be understood in the same language of renewable motifs and filtered exits that already appears in the particle-lifetime program.

What remains open is substantial but now much better organized. The branching-volume problem, the derivation of $\sigma$-filters from topology and ledger structure, the full electroweak density bridge, and the sharper empirical bite of higher-order persistence ideas all remain targets. But these targets now belong to a visible architecture rather than a fog of ambition. That is already a significant gain.

If Relational Actualism is right, the universe is not a static inventory of things arranged in spacetime and governed by externally imposed laws. It is a self-writing growth process whose stable structures survive because graph dynamics filters their destruction. This paper has attempted to show what such a world should look like when it meets experiment.

\begin{thebibliography}{99}

\bibitem{BenincasaDowker2010}
D.~M.~T.~Benincasa and F.~Dowker,
``The scalar curvature of a causal set,''
\emph{Phys. Rev. Lett.} \textbf{104} (2010) 181301.

\bibitem{DowkerGlaser2013}
F.~Dowker and L.~Glaser,
``Causal set d'Alembertians for various dimensions,''
\emph{Class. Quantum Grav.} \textbf{30} (2013) 195016.

\bibitem{RAKernel}
J.~Sandeman,
``Relational Actualism: A process ontology for physics,''
companion theory-definition paper.

\bibitem{RABridge}
J.~Sandeman,
``From the growing ledger to known physics: gravity, matter, and coupling structure in Relational Actualism,''
companion bridge paper.

\bibitem{RAActualizationAudit}
J.~Sandeman,
``The actualization criterion in Relational Actualism,''
internal audit note.

\bibitem{RAd4Closure}
J.~Sandeman,
``d=4 closure in Relational Actualism,''
internal technical note.

\bibitem{RAHadronicDensity}
J.~Sandeman,
``The hadronic density closure in Relational Actualism,''
internal technical note.

\bibitem{RAMotifRenewal}
J.~Sandeman,
``Motif renewal and topology-space decay in Relational Actualism,''
internal technical note.

\bibitem{RASelectionRules}
J.~Sandeman,
``Selection rules and pathwise exit in Relational Actualism,''
internal technical note.

\bibitem{RAGRBridge}
J.~Sandeman,
``The GR bridge in Relational Actualism,''
internal audit note.

\bibitem{RATopoAudit}
J.~Sandeman,
``Particle-topology audit,''
internal audit note.

\bibitem{RAGS02}
J.~Sandeman,
``GS02: Gauge structure in Relational Actualism,''
internal project note.

\bibitem{RAIC30}
J.~Sandeman,
``IC30 and the electromagnetic coupling,''
internal project note.

\bibitem{RACausalFirewall}
J.~Sandeman,
``The causal firewall and higher-order persistence,''
internal project note.

\end{thebibliography}

\end{document}
